\chapter{JavaScript avanzato}

Nel capitolo precedente abbiamo visto i concetti chiave del linguaggio:
variabili, funzioni, scope (e hoisting), oggetti e costruttori. Qui
proseguiamo con altri concetti fondamentali: i \textbf{prototipi} e
l'ereditarietà, le \textbf{strutture dati} (array, mappe, insiemi), le
\textbf{classi}, la gestione degli \textbf{errori} e i \textbf{moduli}.

\section{Prototipi ed ereditarietà}

L'ereditarietà è un aspetto essenziale della programmazione orientata agli
oggetti. JavaScript la realizza in un modo particolare, detto
\term{ereditarietà prototipale}.

\dfn{Prototipo}{
  Ogni oggetto ha una speciale proprietà nascosta chiamata
  \term{[[Prototype]]}, che vale \code{null} oppure fa riferimento a un altro
  oggetto. Quando accediamo a una proprietà di un oggetto e questa manca,
  JavaScript la cerca risalendo la \term{catena dei prototipi}.}

La ricerca di una proprietà risale la catena un anello alla volta, finché non
la trova o finché non arriva a \code{null}.

\begin{esempio}{Impostare un prototipo}
\begin{affianco}[0.6][c]
\begin{lstlisting}[style=js, name=prototipi.js]
let pet = {
  legs: 4,
  greet() { console.log("..."); }
}

let cat = {
  __proto__: pet          // imposta il prototipo
}

cat.greet();              // "..."
console.log(cat.legs);    // 4
console.log(cat.__proto__);  // Object { legs: 4, greet: greet() }
\end{lstlisting}
\risultato
\begin{immagine}[\code{cat} non ha \code{legs} né \code{greet}:\\
                 li trova risalendo la catena dei prototipi.]
\begin{tikzpicture}[
    pr/.style={-{Stealth[length=5pt,width=4pt]}, line width=0.85pt,
      draw=accent!70!black}]
  \mchiave{1.0cm}
  \oggetto{c}{(0,0)}{2.6cm}{cat}{}
  \oggetto{p}{(0,-1.65)}{2.6cm}{pet}{legs/4, greet/function}
  \oggetto[fill=neutral!12!white]{op}{(0,-3.75)}{2.6cm}{Object.prototype}
    {toString/function, \ldots/}
  \node[mcod, text=neutral!35!black] (n) at (0,-5.55) {null};
  \draw[pr] (c-box.south) -- node[wtlab, right, xshift=2pt]{\code{[[Prototype]]}} (p.north);
  \draw[pr] (p-box.south) -- node[wtlab, right, xshift=2pt]{\code{[[Prototype]]}} (op.north);
  \draw[pr] (op-box.south) -- node[wtlab, right, xshift=2pt]{\code{[[Prototype]]}} (n.north);
  \draw[mcerca] ([xshift=-4pt]c-1.west) to[out=180, in=180]
    node[mnota, left]{\code{cat.legs}?} ([xshift=-4pt]p-1.west);
\end{tikzpicture}
\end{immagine}
\end{affianco}

\code{\_\_proto\_\_} è un getter/setter per la vera proprietà
\code{[[Prototype]]}. Nel JavaScript moderno si preferiscono
\code{Object.getPrototypeOf()} e \code{Object.setPrototypeOf()}.
\end{esempio}

\subsection{I prototipi si usano solo in lettura}

Questo è il punto che genera più confusione: i prototipi intervengono
\textbf{solo} quando si legge una proprietà. Le operazioni di scrittura e di
cancellazione agiscono \textbf{direttamente sull'oggetto}.

\begin{esempio}{Scrittura e prototipo}
\begin{affianco}[0.55][c]
\begin{lstlisting}[style=js, name=scrittura-proto.js]
let pet = { legs: 4 }

let cat   = { name: "Garfield" }
let snake = { name: "Salazar" }

Object.setPrototypeOf(cat, pet);
Object.setPrototypeOf(snake, pet);

snake.legs = 0;   // NON tocca pet: crea legs su snake

console.log(`Cat legs: ${cat.legs}`);      // 4 (dal prototipo)
console.log(`Snake legs: ${snake.legs}`);  // 0 (propria)
console.log(`Pet legs: ${pet.legs}`);      // 4 (intatta)
\end{lstlisting}
\risultato
\begin{immagine}[La scrittura crea \code{legs} su \code{snake}:\\
                 il prototipo resta com'era.]
\begin{tikzpicture}[
    pr/.style={-{Stealth[length=5pt,width=4pt]}, line width=0.85pt,
      draw=accent!70!black}]
  \mchiave{1.0cm}
  \oggetto{p}{(1.7,0)}{2.3cm}{pet}{legs/4}
  \oggetto{c}{(0,-2.2)}{2.6cm}{cat}{name/"Garfield"}
  \oggetto{s}{(3.4,-2.2)}{2.6cm}{snake}{name/"Salazar", legs/0}
  \node[draw=accent!60!black, line width=1.1pt, fit=(s-2), inner sep=0pt] {};
  \node[mnota, anchor=west, text=accent!60!black] at ([xshift=3pt]s-2.east)
    {propria};
  \draw[pr] (c.north) -- ([xshift=0.35cm]p-box.south west);
  \draw[pr] (s.north) -- ([xshift=-0.35cm]p-box.south east);
  \node[wtlab] at (1.7,-1.3) {\code{[[Prototype]]}};
\end{tikzpicture}
\end{immagine}
\end{affianco}

Assegnando \code{snake.legs = 0} non si è toccato \code{pet}: si è creata una
proprietà \textbf{propria} su \code{snake}, che da quel momento
\guillemotleft oscura\guillemotright\ quella del prototipo.
\end{esempio}

\subsection{Costruttori e prototipi}

Possiamo creare oggetti con le funzioni costruttore, come \code{new Pet()}.
Quando \code{Pet.prototype} è un oggetto, l'operatore \code{new} lo usa per
impostare il \code{[[Prototype]]} dell'oggetto appena creato.

\begin{affianco}[0.53][c]
\begin{lstlisting}[style=js, name=costruttore-proto.js]
let pet = { eats: true };

function Duck(name) {
  this.name = name;
}

Duck.prototype = pet;

let chuck = new Duck("Chuck");   // chuck.__proto__ == pet
console.log(chuck.eats);         // true
\end{lstlisting}
\risultato
\begin{immagine}[\code{new} collega \code{chuck} all'oggetto\\
                 contenuto in \code{Duck.prototype}.]
\begin{tikzpicture}[
    pr/.style={-{Stealth[length=5pt,width=4pt]}, line width=0.85pt,
      draw=accent!70!black}]
  \mchiave{0.95cm}
  \node[mfun] (d) at (0,0) {Duck};
  \oggetto{p}{(3.6,0)}{2.0cm}{pet}{eats/true}
  \oggetto{c}{(3.6,-2.0)}{2.4cm}{chuck}{name/"Chuck"}
  \draw[pr] (d.east) -- node[wtlab, above]{\code{prototype}} (p-box.west |- d.east);
  \draw[pr] (c.north) -- node[wtlab, right, xshift=2pt]{\code{[[Prototype]]}} (p-box.south);
  \draw[mrif, dashed] (d.south) |- node[wtlab, pos=0.25, left]{\code{new}} (c-box.west);
\end{tikzpicture}
\end{immagine}
\end{affianco}

E se alla funzione costruttore non assegniamo nessun \code{prototype}? In quel
caso viene usato un \term{prototipo predefinito}: un oggetto con una sola
proprietà \code{constructor}, che punta indietro alla funzione costruttore. Ogni
funzione ordinaria lo riceve automaticamente nel momento in cui viene creata, ed
è come se JavaScript eseguisse per noi:

\begin{lstlisting}[style=js]
Pet.prototype = { constructor: Pet };   // fatto in automatico
\end{lstlisting}

L'unica differenza rispetto a questa riga scritta a mano è che, nel prototipo
predefinito, \code{constructor} non è \textbf{enumerabile}: non compare in un
ciclo \code{for..in} né in \code{Object.keys}.

\warning{Sostituire il prototipo fa perdere constructor}{
  Nell'esempio di \code{Duck} il prototipo predefinito è stato
  \textbf{sostituito} con \code{pet}, che non ha una proprietà
  \code{constructor}. La ricerca allora risale fino a \code{Object.prototype}:
  \code{chuck.constructor} vale \code{Object}, non \code{Duck}. Se serve, va
  rimessa a mano: \code{Duck.prototype.constructor = Duck}. Oppure, più
  semplicemente, si \textbf{aggiungono} proprietà al prototipo esistente invece
  di sostituirlo.}

Il codice a sinistra crea esattamente gli oggetti dello schema a destra: una
funzione costruttore \code{Pet} senza un \code{prototype} assegnato a mano, e due
oggetti creati con \code{new}. I controlli in fondo verificano, una per una, le
frecce del disegno.

\begin{affianco}[0.55][c]
\codicepiccolo
\begin{lstlisting}[style=js, name=costruttore-predefinito.js]
function Pet(name, age) {
  this.name = name;
  this.age = age;
}
// nessuna assegnazione a Pet.prototype: si usa quello predefinito

let chuck = new Pet("Chuck", 3);
let rick  = new Pet("Rick", 7);

// le frecce dello schema
Pet.prototype.constructor === Pet;                          // true
Object.getPrototypeOf(chuck) === Pet.prototype;             // true
Object.getPrototypeOf(rick)  === Pet.prototype;             // true
Object.getPrototypeOf(Pet.prototype) === Object.prototype;  // true

// constructor non e' una proprieta' di chuck: arriva dal prototipo
chuck.hasOwnProperty("constructor");   // false
chuck.constructor === Pet;             // true
chuck.toString();     // "[object Object]", da Object.prototype
\end{lstlisting}
\risultato
\begin{immagine}[Costruttore, prototipo predefinito e oggetti creati.]
\begin{tikzpicture}[
    pr/.style={-{Stealth[length=5pt,width=4pt]}, line width=0.85pt,
      draw=accent!70!black},
    a/.style={-{Stealth[length=5pt,width=4pt]}, line width=0.85pt,
      draw=neutral!45!black}]
  \mchiave{1.0cm}
  \oggetto[fill=neutral!12!white]{op}{(2.2,2.9)}{2.8cm}{Object.prototype}
    {toString/function}
  \oggetto[fill=accent!14!white]{pp}{(2.2,0.9)}{2.8cm}{Pet.prototype}
    {constructor/Pet}
  \node[mfun] (fn) at (-1.35,0.67) {Pet};
  \oggetto{c}{(0.85,-1.2)}{2.3cm}{chuck}{name/"Chuck", age/3}
  \oggetto{r}{(3.55,-1.2)}{2.3cm}{rick}{name/"Rick", age/7}
  \draw[pr] (pp.north) -- node[wtlab, right, xshift=2pt]{\code{[[Prototype]]}}
    (op-box.south);
  \draw[a] ([yshift=5pt]fn.east) -- node[wtlab, above]{\code{prototype}}
    ([yshift=5pt]fn.east -| pp-box.west);
  \draw[a] ([yshift=-5pt]fn.east -| pp-box.west) -- node[wtlab, below]{\code{constructor}}
    ([yshift=-5pt]fn.east);
  \draw[pr] (c.north) -- ([xshift=-0.5cm]pp-box.south);
  \draw[pr] (r.north) -- ([xshift=0.5cm]pp-box.south);
\end{tikzpicture}
\end{immagine}
\end{affianco}

Questo meccanismo ci permette, dato un oggetto, di risalire al costruttore che
lo ha creato, e di usarlo per crearne uno nuovo dello stesso tipo:

\begin{lstlisting}[style=js]
let b = new chuck.constructor('Bob', 1);   // come new Pet('Bob', 1)
\end{lstlisting}

\begin{warningbox}{{Senza new non funziona}}
  Le slide scrivono \code{chuck.constructor('Bob', 1)}, senza \code{new}. Così
  \code{Pet} viene chiamata come un normale \textbf{metodo} di \code{chuck}:
  \code{this} vale \code{chuck}, la funzione \textbf{sovrascrive} nome ed età di
  \code{chuck} e restituisce \code{undefined}.

\begin{lstlisting}[style=js]
let b = chuck.constructor('Bob', 1);
console.log(b);            // undefined
console.log(chuck.name);   // "Bob" (!)
\end{lstlisting}
\end{warningbox}

\subsection{Due modi per aggiungere metodi}

\begin{affianco}[0.55][c]
\begin{lstlisting}[style=js, name=metodi-proto.js]
function describe() {
  console.log(`I'm ${this.name} and my age is ${this.age}`);
}

function Pet(name, age) {
  this.name = name;
  this.age  = age;
  this.describe = describe;   // (1) su OGNI oggetto creato
}

Pet.prototype.describe = describe;   // (2) uno solo, condiviso

let chuck = new Pet('Chuck', 3);
let rick  = new Pet('Rick', 7);

chuck.describe();
rick.describe();
\end{lstlisting}
\risultato
\begin{immagine}[Con (1) ogni oggetto ha il suo \code{describe};\\
                 con (2) basta quello sul prototipo.]
\begin{tikzpicture}[
    pr/.style={-{Stealth[length=5pt,width=4pt]}, line width=0.85pt,
      draw=accent!70!black}]
  \mchiave{1.15cm}
  \oggetto[fill=accent!14!white]{pp}{(1.6,1.6)}{2.6cm}{Pet.prototype}
    {describe/}
  \oggetto{c}{(0,-0.5)}{2.4cm}{chuck}{name/"Chuck", age/3, describe/}
  \oggetto{r}{(3.2,-0.5)}{2.4cm}{rick}{name/"Rick", age/7, describe/}
  \node[mfun] (f) at (1.6,-3.2) {function describe};
  \draw[pr] (c.north) -- ([xshift=-0.45cm]pp-box.south);
  \draw[pr] (r.north) -- ([xshift=0.45cm]pp-box.south);
  \node (d1) at ([xshift=-9pt]c-3.east) {};
  \node (d2) at ([xshift=-9pt]r-3.east) {};
  \node (dp) at ([xshift=-9pt]pp-1.east) {};
  \riferimento{d1}{to[out=-90, in=150] (f.north west)}
  \riferimento{d2}{to[out=-90, in=30] (f.north east)}
  \riferimento{dp}{to[out=0, in=0, looseness=1.6] (f.east)}
  \node[mnota, text=accent!60!black] at (-1.65,-1.4) {(1)};
  \node[mnota, text=accent!60!black] at (4.9,1.6) {(2)};
\end{tikzpicture}
\end{immagine}
\end{affianco}

\info{Quale dei due conviene}{
  Con la soluzione (1) ogni oggetto riceve una \textbf{propria copia} del
  riferimento al metodo; con la (2) esiste un solo metodo, condiviso da tutti
  gli oggetti tramite il prototipo. Con mille oggetti la differenza in memoria è
  concreta, ed è il motivo per cui i metodi vanno di norma sul prototipo. È
  esattamente ciò che fanno le classi, che vedremo fra poco.}

\section{Array}

Gli \term{array} permettono di memorizzare sequenze ordinate di valori. Si
dichiarano con la sintassi letterale a parentesi quadre o con il costruttore
\code{Array}.

\begin{affianco}[0.55][c]
\begin{lstlisting}[style=js, name=array.js]
// sintassi letterale
let a = ["HTML", "CSS", "JS"];
console.log(a);   // Array(3) [ "HTML", "CSS", "JS" ]

// sintassi con costruttore
let b = new Array("HTML", "CSS", "JS");
console.log(b);   // Array(3) [ "HTML", "CSS", "JS" ]
\end{lstlisting}
\risultato
\begin{immagine}[Tre caselle, numerate a partire da 0.]
\begin{tikzpicture}
  \node[mnome, anchor=east] at (-0.1,0) {a};
  \celle{a}{(0,0)}{"HTML", "CSS", "JS"}
  \node[mnota, anchor=west] at ([xshift=8pt]a-box.east) {\code{length}: 3};
\end{tikzpicture}
\end{immagine}
\end{affianco}

\subsection{Indicizzazione}

Gli array sono indicizzati \textbf{a partire da 0}. I singoli valori si leggono
e si scrivono con la notazione a parentesi quadre. La proprietà \code{length}
contiene l'\textbf{indice massimo più uno}.

\begin{affianco}[0.55][c]
\begin{lstlisting}[style=js, name=indici.js]
let a = ["HTML", "CSS", "JS"];

a[2] = "JavaScript";
a[3] = "React";

console.log(a[0]);       // HTML
console.log(a[2]);       // JavaScript
console.log(a[3]);       // React
console.log(a.length);   // 4
\end{lstlisting}
\risultato
\begin{immagine}[Dopo le due assegnazioni: \code{a[2]} è stata\\
                 sovrascritta, \code{a[3]} è nuova.]
\begin{tikzpicture}
  \node[mnome, anchor=east] at (-0.1,0) {a};
  \celle{a}{(0,0)}{"HTML", "CSS", "JavaScript", "React"}
  \node[draw=accent!60!black, line width=1.1pt, fit=(a-2)(a-3), inner sep=0pt] {};
  \node[mnota, anchor=west] at ([xshift=8pt]a-box.east) {\code{length}: 4};
\end{tikzpicture}
\end{immagine}
\end{affianco}

\subsection{La proprietà length}

Si sarà notata la definizione insolita di \code{length}: non è il
\textbf{conteggio} dei valori nell'array.

\begin{esempio}{length non conta gli elementi}
\begin{affianco}[0.5][c]
\begin{lstlisting}[style=js, name=length.js]
let a = [1, 2, 3, 4, 5];
a[999] = 998;

console.log(a.length);   // 1000 (!)
console.log(a[5]);       // undefined
console.log(a[999]);     // 998
console.log(a[2000]);    // undefined
\end{lstlisting}
\risultato
\begin{immagine}[\code{length} vale l'indice più alto più uno,\\
                 anche se in mezzo ci sono 994 caselle vuote.]
\begin{tikzpicture}
  \noindici
  \node[mnome, anchor=east] at (-0.1,0) {a};
  \celle{a}{(0,0)}{1, 2, 3, 4, 5, \mvuota\ \ldots\ \mvuota, 998}
  \node[mcella, fill=neutral!10!white, draw=neutral!50!black, fit=(a-5),
        inner sep=0pt] {};
  \node at (a-5) {\Lato\scriptsize\itshape\color{neutral!45!black}994 caselle vuote};
  \foreach \i in {0,...,4} \node[mindice, anchor=north] at ([yshift=-1pt]a-\i.south) {\i};
  \node[mindice, anchor=north] at ([yshift=-1pt]a-5.south) {5 \ldots\ 998};
  \node[mindice, anchor=north] at ([yshift=-1pt]a-6.south) {999};
  \node[mnota, anchor=north] at ([yshift=-14pt]a-box.south) {\code{length}: 1000};
\end{tikzpicture}
\end{immagine}
\end{affianco}

Un'altra cosa interessante è che \code{length} è \textbf{scrivibile}:

\begin{affianco}[0.5][c]
\begin{lstlisting}[style=js, name=length2.js]
let a = [1, 2, 3, 4, 5];

a.length = 3;
console.log(a);    // [1, 2, 3]

a.length = 5;
console.log(a);    // [1, 2, 3, <2 empty slots>]

a.length = 0;      // svuota l'array
\end{lstlisting}
\risultato
\begin{immagine}
\begin{tikzpicture}
  \node[mcod, anchor=east] at (-0.2,0) {length = 3};
  \celle{a}{(0,0)}{1, 2, 3}
  \node[mcod, anchor=east] at (-0.2,-0.95) {length = 5};
  \celle{b}{(0,-0.95)}{1, 2, 3, \mvuota, \mvuota}
  \node[mcod, anchor=east] at (-0.2,-1.9) {length = 0};
  \node[mnota, anchor=west] at (0,-1.9) {(nessuna casella)};
\end{tikzpicture}
\end{immagine}
\end{affianco}

Assegnare a \code{length} un valore più piccolo è il modo più rapido (e più
sorprendente) di troncare un array.
\end{esempio}

\subsection{Tipi eterogenei}

Un array può contenere dati di tipo diverso, funzioni comprese.

\begin{affianco}[0.55][c]
\begin{lstlisting}[style=js, name=eterogeneo.js]
let a = [
  "JavaScript",
  { name: "John", job: "Dev" },
  12,
  function (name) { console.log(`Hello ${name}`); }
]

console.log(a[1].name);        // John
a[3]("Web Technologies");      // Hello Web Technologies
\end{lstlisting}
\risultato
\begin{immagine}[Una casella può contenere qualunque valore,\\
                 anche un riferimento a un oggetto o a una funzione.]
\begin{tikzpicture}
  \mchiave{0.75cm}
  \node[mnome, anchor=east] at (-0.1,0) {a};
  \celle{a}{(0,0)}{"JavaScript", \phantom{xx}, 12, \phantom{xx}}
  \oggetto{o}{(1.6,-1.7)}{2.2cm}{Object}{name/"John", job/"Dev"}
  \node[mfun] (f) at (4.3,-1.95) {function (name)};
  \riferimento{a-1}{-- (a-1 |- o.north)}
  \riferimento{a-3}{-- (a-3 |- f.north)}
\end{tikzpicture}
\end{immagine}
\end{affianco}

\subsection{Metodi per aggiungere e rimuovere}

\begin{center}
\renewcommand{\arraystretch}{1.35}
\begin{tabular}{@{}l@{\hspace{1.2cm}}l@{}}
  \textbf{Metodo} & \textbf{Effetto} \\
  \code{push()}    & aggiunge un elemento \textbf{in fondo} \\
  \code{pop()}     & preleva un elemento \textbf{dal fondo} \\
  \code{unshift()} & aggiunge un elemento \textbf{all'inizio} \\
  \code{shift()}   & preleva un elemento \textbf{dall'inizio} \\
\end{tabular}
\end{center}

\begin{affianco}[0.45][c]
\begin{lstlisting}[style=js, name=metodi-array.js]
let a = [1];      // [1]
a.push(2);        // [1, 2]
a.unshift(0);     // [0, 1, 2]

let x = a.pop();     // [0, 1]
let y = a.shift();   // [1]

console.log(x);   // 2
console.log(y);   // 0
\end{lstlisting}
\risultato
\begin{immagine}[I quattro metodi fondamentali (in viola\\
                 l'elemento aggiunto o rimosso).]
\begin{tikzpicture}[font=\FiraCodeNL\footnotesize,
      c/.style={draw=primary!60!black, fill=primary!10!white, line width=0.7pt,
                minimum size=0.55cm},
      n/.style={c, draw=accent!60!black, fill=accent!16!white},
      a/.style={-{Stealth[length=5pt,width=4pt]}, line width=0.8pt,
                draw=neutral!45!black}]

    \foreach \op/\i in {push/0, unshift/1, pop/2, shift/3} {
      \node[font=\FiraCodeNL\scriptsize, anchor=east, text=neutral!25!black]
        at (-0.25,-\i*1.05) {.\op()};
    }

    % push
    \begin{scope}[yshift=0cm]
      \node[c] (p1) at (0.3,0) {A}; \node[c] (p2) at (0.85,0) {B};
      \draw[a] (1.5,0) -- (2.4,0);
      \node[c] at (2.9,0) {A}; \node[c] at (3.45,0) {B}; \node[n] at (4.0,0) {C};
    \end{scope}
    % unshift
    \begin{scope}[yshift=-1.05cm]
      \node[c] at (0.3,0) {A}; \node[c] at (0.85,0) {B};
      \draw[a] (1.5,0) -- (2.4,0);
      \node[n] at (2.9,0) {C}; \node[c] at (3.45,0) {A}; \node[c] at (4.0,0) {B};
    \end{scope}
    % pop
    \begin{scope}[yshift=-2.1cm]
      \node[c] at (0.3,0) {A}; \node[c] at (0.85,0) {B}; \node[n] at (1.4,0) {C};
      \draw[a] (2.0,0) -- (2.9,0);
      \node[c] at (3.4,0) {A}; \node[c] at (3.95,0) {B};
    \end{scope}
    % shift
    \begin{scope}[yshift=-3.15cm]
      \node[n] at (0.3,0) {A}; \node[c] at (0.85,0) {B}; \node[c] at (1.4,0) {C};
      \draw[a] (2.0,0) -- (2.9,0);
      \node[c] at (3.4,0) {B}; \node[c] at (3.95,0) {C};
    \end{scope}
  \end{tikzpicture}
\end{immagine}
\end{affianco}

\info{Tutti i metodi, visualizzati}{
  Gli array hanno molti altri metodi (\code{slice}, \code{splice},
  \code{map}, \code{filter}, \code{find}, \code{includes}, \code{sort},
  \ldots). Una visualizzazione di ciascuno, nello stesso stile dello schema, si
  trova su \urlx{https://js-arrays-visualized.com/}.}

\subsection{Iterazione}

Gli array si possono percorrere con un \code{for} sugli indici, con la sintassi
\code{for..of}, oppure con il metodo \code{forEach}.

\begin{affianco}[0.55][c]
\begin{lstlisting}[style=js, name=iterazione.js]
let a = ["a", "b", "c"];
a[4] = "e";              // lascia un buco in posizione 3

for (let i = 0; i < a.length; i++)
  console.log(a[i]);     // a, b, c, undefined, e

for (let item of a)
  console.log(item);     // a, b, c, undefined, e

a.forEach((value, index, array) => {
  console.log(`a[${index}]=${value}`);
});

for (let key in a) {
  console.log(a[key]);   // a, b, c, e   (salta il buco!)
}
\end{lstlisting}
\risultato
\begin{immagine}[\code{for}, \code{for..of} e \code{forEach} passano da\\
                 tutti gli indici; \code{for..in} salta il buco.]
\begin{tikzpicture}[
    si/.style={circle, fill=mok, inner sep=1.6pt},
    lab/.style={mnota, anchor=east}]
  \node[mnome, anchor=east] at (-0.1,0) {a};
  \celle{a}{(0,0)}{"a", "b", "c", \mvuota, "e"}
  \node[lab] at (-0.1,-0.9) {\code{for}, \code{for..of}};
  \node[lab] at (-0.1,-1.75) {\code{for..in}};
  \foreach \i in {0,...,4} {
    \node[si] at (a-\i |- 0,-0.9) {};
  }
  \node[mindice, anchor=north, text=neutral!40!black] at ([yshift=-4pt]a-3 |- 0,-0.9)
    {undefined};
  \foreach \i in {0,1,2,4} {
    \node[si] at (a-\i |- 0,-1.75) {};
  }
  \node at (a-3 |- 0,-1.75) {\merr};
\end{tikzpicture}
\end{immagine}
\end{affianco}

\warning{Non usare for..in sugli array}{
  Poiché gli array sono oggetti, \code{for..in} funziona, ma è una cattiva idea
  per due motivi. È da 10 a 100 volte \textbf{più lento}, e itera su
  \textbf{tutte} le proprietà enumerabili, comprese quelle che non sono valori
  dell'array. Come si vede nell'esempio, salta anche i buchi, comportandosi in
  modo diverso da \code{for..of}. Per gli array si usa \code{for..of} o
  \code{forEach}; \code{for..in} è per gli oggetti.}

\subsection{Array multidimensionali}

Gli elementi di un array possono essere altri array: è così che si definiscono
le matrici.

\begin{affianco}[0.55][c]
\begin{lstlisting}[style=js, name=matrice.js]
let matrix = [
  [1, 2, 3],
  [4, 5, 6],
  [7, 8, 9],
]

console.log(matrix[0][0]);   // 1
console.log(matrix[1][1]);   // 5
console.log(matrix[2][2]);   // 9
\end{lstlisting}
\risultato
\begin{immagine}[Un array di array: il primo indice sceglie la riga,\\
                 il secondo la colonna.]
\begin{tikzpicture}[
    m/.style={mcella, minimum width=0.62cm},
    d/.style={m, fill=accent!14!white, draw=accent!55!black}]
  \foreach \r in {0,1,2} {
    \node[mcod, anchor=east] at (-0.15,-\r*0.62) {matrix[\r]};
    \foreach \c in {0,1,2} {
      \pgfmathtruncatemacro{\v}{\r*3+\c+1}
      \ifnum\r=\c
        \node[d] at (\c*0.62+0.31,-\r*0.62) {\v};
      \else
        \node[m] at (\c*0.62+0.31,-\r*0.62) {\v};
      \fi
    }
  }
  \foreach \c in {0,1,2}
    \node[mindice] at (\c*0.62+0.31,0.45) {[\c]};
\end{tikzpicture}
\end{immagine}
\end{affianco}

\section{Assegnamenti destrutturanti}

Una sintassi speciale che permette di \guillemotleft spacchettare\guillemotright\ un array in più
variabili.

\begin{affianco}[0.56][c]
\begin{lstlisting}[style=js, name=destrutturazione.js]
let [x, y] = ["a", "b", "c"];   // l'eccesso viene scartato
console.log(x);      // a
console.log(y);      // b

let [a, b, ...rest] = [1, 2, 3, 4, 5];   // il resto in rest
console.log(a);      // 1
console.log(b);      // 2
console.log(rest);   // [3, 4, 5]

// la stessa sintassi vale per i parametri di una funzione
function greet(msg, ...names) {
  console.log(`${msg} to ${names}`);
}

greet("Hello", "Ann", "Bob", "Carl");   // Hello to Ann,Bob,Carl
\end{lstlisting}
\risultato
\begin{immagine}[Ogni variabile prende la casella nella sua posizione;\\
                 \code{...rest} raccoglie tutte quelle rimaste.]
\begin{tikzpicture}[
    a/.style={-{Stealth[length=4pt,width=3.5pt]}, line width=0.75pt,
      draw=accent!70!black},
    v/.style={mvar, minimum width=0.62cm}]
  \noindici
  % [x, y] = ["a", "b", "c"]
  \celle{p}{(0,0)}{"a", "b", "c"}
  \node[v] (x) at (p-0 |- 0,-1.05) {"a"};
  \node[mnome, anchor=north] at (x.south) {x};
  \node[v] (y) at (p-1 |- 0,-1.05) {"b"};
  \node[mnome, anchor=north] at (y.south) {y};
  \draw[a] (p-0) -- (x); \draw[a] (p-1) -- (y);
  \node[mnota, text=merr, anchor=north] at ([yshift=-3pt]p-2.south) {scartato};
  % [a, b, ...rest] = [1, 2, 3, 4, 5]
  \celle{q}{(2.9,0)}{1, 2, 3, 4, 5}
  \node[v] (va) at (q-0 |- 0,-1.05) {1};
  \node[mnome, anchor=north] at (va.south) {a};
  \node[v] (vb) at (q-1 |- 0,-1.05) {2};
  \node[mnome, anchor=north] at (vb.south) {b};
  \draw[a] (q-0) -- (va); \draw[a] (q-1) -- (vb);
  \node[mvar, minimum width=1.6cm] (vr) at (q-3 |- 0,-1.05) {[3, 4, 5]};
  \node[mnome, anchor=north] at (vr.south) {rest};
  \draw[wtbrace] ([yshift=-1pt]q-4.south east) -- ([yshift=-1pt]q-2.south west);
  \draw[a] ([yshift=-5pt]q-3.south) -- (vr);
\end{tikzpicture}
\end{immagine}
\end{affianco}

\section{Iterabili}

Il ciclo \code{for..of} funziona con gli \term{oggetti iterabili}. Gli array
sono iterabili, e lo sono anche le stringhe.

\begin{affianco}[0.5][c]
\begin{lstlisting}[style=js, name=iterabili.js]
let string = "Web Technologies!";
for (let char of string) {
  console.log(char);   // W, e, b, , T, e, ...
}
\end{lstlisting}
\risultato
\begin{immagine}[Una stringa è iterabile: \code{for..of} la visita\\
                 un carattere alla volta.]
\begin{tikzpicture}[
    a/.style={-{Stealth[length=4pt,width=3.5pt]}, line width=0.75pt,
      draw=accent!70!black}]
  \noindici
  \tikzset{mcella/.append style={minimum width=0.5cm}}
  \celle{s}{(0,0)}{W, e, b, \textvisiblespace, T, e, c, h, \ldots}
  \foreach \i/\j in {0/1, 1/2, 2/3, 3/4, 4/5} {
    \draw[a] ([yshift=-3pt]s-\i.south) to[out=-60, in=-120] ([yshift=-3pt]s-\j.south);
  }
  \node[mnota, anchor=north west] at ([yshift=-14pt]s-0.south west) {\code{char} assume un valore per giro};
\end{tikzpicture}
\end{immagine}
\end{affianco}

E se avessimo un oggetto nostro che rappresenta una collezione di valori e
volessimo usarlo in un \code{for..of}?

\begin{lstlisting}[style=js, name=range.js]
// range rappresenta un insieme di interi: da min fino a max, a passi di step
let range = { min: 1, max: 10, step: 2 }

for (let value of range) {   // TypeError: range is not iterable
  console.log(value);
}
\end{lstlisting}

Per rendere un oggetto iterabile bisogna implementare il metodo speciale
\code{Symbol.iterator}.

\subsection{Il protocollo di iterazione}

\begin{itemize}
  \item Quando un ciclo \code{for..of} inizia, invoca \textbf{una sola volta} il
        metodo \code{Symbol.iterator} sull'oggetto (e solleva un
        \code{TypeError} se il metodo non esiste).
  \item Da lì in avanti il ciclo lavora soltanto con l'oggetto
        (\term{iteratore}) restituito da quella chiamata.
  \item Quando serve il valore successivo, il ciclo invoca il metodo
        \code{next()} sull'iteratore.
  \item Il risultato di \code{next()} è un oggetto della forma
        \code{\{done: boolean, value: any\}}, dove \code{done: true} significa
        che l'iterazione è finita; altrimenti \code{value} è il valore
        successivo.
\end{itemize}

\begin{esempio}{Implementare Symbol.iterator}
\begin{affianco}[0.5][c]
\begin{lstlisting}[style=js, name=symbol-iterator.js]
let range = { min: 1, max: 10, step: 3 }

range[Symbol.iterator] = function () {
  return {
    current: this.min,
    max: this.max,
    step: this.step,

    next() {
      let n = { done: false, value: this.current };
      if (n.value > this.max)
        n.done = true;
      this.current += this.step;
      return n;
    }
  }
}

for (let value of range)
  console.log(value);   // 1, 4, 7, 10
\end{lstlisting}
\risultato
\begin{immagine}[Il dialogo fra il ciclo e l'iteratore: una chiamata\\
                 a \code{Symbol.iterator}, poi \code{next()} a ogni giro.]
\begin{tikzpicture}[
    att/.style={rounded corners=3pt, draw=primary!65!black, fill=primary!12!white,
      line width=0.7pt, font=\Lato\scriptsize\bfseries, minimum width=2.2cm,
      minimum height=0.55cm},
    vita/.style={dashed, draw=neutral!55!white, line width=0.7pt},
    chiama/.style={-{Stealth[length=5pt,width=4pt]}, line width=0.8pt,
      draw=primary!70!black},
    torna/.style={-{Stealth[length=5pt,width=4pt]}, line width=0.7pt,
      draw=neutral!50!black, dashed},
    msg/.style={mcod, above, inner sep=1.5pt}]
  \node[att] (f) at (0,0) {for..of};
  \node[att] (r) at (5.2,0) {range};
  \draw[vita] (f.south) -- (0,-6.5);
  \draw[vita] (r.south) -- (5.2,-6.5);
  \draw[chiama] (0,-0.75) -- node[msg]{range[Symbol.iterator]()} (5.2,-0.75);
  \node[att, fill=accent!12!white, draw=accent!55!black] (it) at (5.2,-1.55) {iteratore};
  \draw[torna] (it.west) -- node[msg]{iteratore} (0,-1.55);
  \foreach \y/\v in {-2.35/{\{done: false, value: 1\}}, -3.35/{\{done: false, value: 4\}},
                     -4.35/{\ldots\ value: 7, poi 10 \ldots}, -5.35/{\{done: true\}}} {
    \draw[chiama] (0,\y) -- node[msg]{next()} (5.2,\y);
    \draw[torna] (5.2,\y-0.45) -- node[msg]{\v} (0,\y-0.45);
  }
  \node[mnota, anchor=west, text=mok] at (0.1,-6.2) {il ciclo termina};
\end{tikzpicture}
\end{immagine}
\end{affianco}

È anche possibile lavorare direttamente con l'iteratore. Il ciclo seguente è
equivalente al \code{for..of} qui sopra:

\begin{lstlisting}[style=js, name=iteratore-esplicito.js]
let iterator = range[Symbol.iterator]();

while (true) {
  let result = iterator.next();
  if (result.done)
    break;
  console.log(result.value);
}
\end{lstlisting}
\end{esempio}

\section{Mappe}

Le \term{mappe} sono collezioni di coppie chiave-valore. Sono quindi come gli
oggetti? Non proprio: le mappe ammettono chiavi di \textbf{qualunque tipo}.

\begin{center}
\renewcommand{\arraystretch}{1.35}
\begin{tabular}{@{}l@{\hspace{1.0cm}}>{\raggedright\arraybackslash}p{8.4cm}@{}}
  \textbf{Metodo} & \textbf{Descrizione} \\
  \code{new Map()}           & crea la mappa \\
  \code{map.set(key, value)} & memorizza il valore alla chiave \\
  \code{map.get(key)}        & restituisce il valore, o \code{undefined} \\
  \code{map.has(key)}        & \code{true} se la chiave esiste \\
  \code{map.delete(key)}     & rimuove la coppia chiave-valore \\
  \code{map.clear()}         & svuota la mappa \\
  \code{map.size}            & il numero di elementi \\
\end{tabular}
\end{center}

\begin{esempio}{Mappe contro oggetti}
\begin{affianco}[0.55][c]
\codicepiccolo
\begin{lstlisting}[style=js, name=mappe.js]
let map = new Map();
let o = {};

map.set(1, "Num");      map.set("1", "Str");
map.set(true, "Bool");  map.set("true", "Str");

console.log(map);
// Map { 1 -> "Num", "1" -> "Str", true -> "Bool", "true" -> "Str" }
console.log(map.size);   // 4

o[1] = "Num";     o["1"] = "String";
o[true] = "Bool"; o["true"] = "String";

console.log(o);   // Object { 1: "String", true: "String" }
\end{lstlisting}
\risultato
\begin{immagine}[Nella mappa le chiavi conservano il tipo; nell'oggetto\\
                 diventano stringhe e due coppie si sovrascrivono.]
\begin{tikzpicture}
  \mchiave{1.05cm}
  \oggetto[fill=accent!14!white]{m}{(0,0)}{2.6cm}{Map (4 chiavi)}
    {1/"Num", "1"/"Str", true/"Bool", "true"/"Str"}
  \oggetto{o}{(3.4,0)}{2.6cm}{Object (2 chiavi)}
    {"1"/"String", "true"/"String"}
  \node[mnota, anchor=north west, align=left] at ([yshift=-4pt]o-box.south west)
    {\code{1} e \code{"1"} sono\\ la stessa chiave \code{"1"}};
\end{tikzpicture}
\end{immagine}
\end{affianco}

La differenza è netta: nell'oggetto le chiavi vengono \textbf{convertite in
stringa}, quindi \code{1} e \code{"1"} sono la stessa chiave e i valori si
sovrascrivono; restano due proprietà su quattro. Nella mappa le quattro chiavi
sono distinte.
\end{esempio}

\subsection{Iterare su una mappa}

\begin{lstlisting}[style=js, name=iterazione-mappe.js]
let map = new Map();
map.set("Hello", "JS"); map.set(1, true); map.set(false, 42);

for (let key of map.keys()) {          // iterabile sulle chiavi
  console.log(key);                    // Hello, 1, false
}

for (let value of map.values()) {      // iterabile sui valori
  console.log(value);                  // JS, true, 42
}

for (let [key, value] of map.entries()) {   // iterabile sulle coppie
  console.log(`${key}: ${value}`);
}

for (let [key, value] of map) {        // equivalente al precedente
  console.log(`${key}: ${value}`);
}
\end{lstlisting}

\section{Insiemi}

Gli \term{insiemi} (\textit{Set}) sono una struttura dati per collezioni di
valori \textbf{senza ripetizioni}.

\begin{center}
\renewcommand{\arraystretch}{1.35}
\begin{tabular}{@{}l@{\hspace{1.0cm}}>{\raggedright\arraybackslash}p{8.0cm}@{}}
  \textbf{Metodo} & \textbf{Descrizione} \\
  \code{new Set([iterable])} & crea l'insieme; se riceve un iterabile, ne copia
                               i valori \\
  \code{set.add(value)}      & memorizza il valore, restituisce l'insieme \\
  \code{set.delete(value)}   & elimina il valore; \code{true} se esisteva \\
  \code{set.has(value)}      & \code{true} se il valore è presente \\
  \code{set.clear()}         & svuota l'insieme \\
  \code{set.size}            & il numero di elementi \\
\end{tabular}
\end{center}

\begin{affianco}[0.55][c]
\begin{lstlisting}[style=js, name=insiemi.js]
let set = new Set();

let eric = { name: "Eric" };
let stan = { name: "Stan" };
let kyle = { name: "Kyle" };

set.add(eric); set.add(stan);
set.add(kyle); set.add(eric);   // gia' presente
set.add(kyle);                  // gia' presente

console.log(set.size);          // 3

for (let user of set) {
  console.log(user.name);       // Eric, Stan, Kyle
}
\end{lstlisting}
\risultato
\begin{immagine}[Aggiungere un valore già presente non cambia nulla.]
\begin{tikzpicture}[
    el/.style={rounded corners=2pt, draw=primary!65!black, fill=white,
      line width=0.7pt, font=\FiraCodeNL\scriptsize, minimum height=0.5cm},
    a/.style={-{Stealth[length=4pt,width=3.5pt]}, line width=0.75pt,
      draw=neutral!45!black}]
  \node[el] (e) at (0,0) {eric};
  \node[el] (s1) at (1.1,0) {stan};
  \node[el] (k) at (2.2,0) {kyle};
  \begin{pgfonlayer}{background}
  \node[draw=accent!60!black, fill=accent!7!white, rounded corners=8pt,
        line width=0.8pt, fit=(e)(s1)(k), inner sep=7pt,
        label={[mnota]above:\code{set} \quad \code{size}: 3}] (set) {};
  \end{pgfonlayer}
  \node[el, draw=neutral!45!white, text=neutral!50!black] (e2) at (4.6,0.4) {eric};
  \node[el, draw=neutral!45!white, text=neutral!50!black] (k2) at (4.6,-0.4) {kyle};
  \draw[a] (e2.west) -- node[fill=white, inner sep=0pt]{\merr} ([yshift=0.25cm]set.east);
  \draw[a] (k2.west) -- node[fill=white, inner sep=0pt]{\merr} ([yshift=-0.25cm]set.east);
  \node[mnota, anchor=west] at ([xshift=4pt]e2.east) {già presente};
  \node[mnota, anchor=west] at ([xshift=4pt]k2.east) {già presente};
\end{tikzpicture}
\end{immagine}
\end{affianco}

\section{Classi}

Conosciamo il concetto di classe dalla programmazione orientata agli oggetti: le
classi sono modelli per creare oggetti. In JavaScript i costruttori e
\code{new} servono già a questo, ma il linguaggio offre anche un costrutto
\code{class} più evoluto.

\begin{affianco}[0.5][c]
\begin{lstlisting}[style=js, name=classi.js]
class Pet {
  constructor(name) { this.name = name }
  method1() { /* ... */ }
  method2() { /* ... */ }
}

let pet = new Pet("Chuck");
\end{lstlisting}
\risultato
\begin{immagine}[Che cosa produce \code{class Pet \{...\}}.]
\begin{tikzpicture}[
    pr/.style={-{Stealth[length=5pt,width=4pt]}, line width=0.85pt,
      draw=accent!70!black},
    a/.style={-{Stealth[length=5pt,width=4pt]}, line width=0.85pt,
      draw=neutral!45!black}]
  \mchiave{1.1cm}
  \node[mfun, align=left] (fn) at (0,0)
    {function Pet(name) \{\\ \ \ this.name = name\\ \}};
  \oggetto[fill=accent!14!white]{pp}{(4.2,0.3)}{2.4cm}{Pet.prototype}
    {method1/function, method2/function, constructor/Pet}
  \draw[a] ([yshift=5pt]fn.east) -- node[wtlab, above]{\code{prototype}}
    ([yshift=5pt]fn.east -| pp-box.west);
  \oggetto{o}{(4.2,-2.0)}{2.4cm}{pet}{name/"Chuck"}
  \draw[pr] (o.north) -- node[wtlab, right, xshift=2pt]{\code{[[Prototype]]}} (pp-box.south);
  \node[mnota, anchor=north] at ([yshift=-3pt]fn.south)
    {il corpo viene dal \code{constructor}};
\end{tikzpicture}
\end{immagine}
\end{affianco}

\code{new Pet()} crea un nuovo oggetto e invoca il metodo \code{constructor}
con gli argomenti dati, che imposta la proprietà \code{name} sull'oggetto. Il
nuovo oggetto viene poi restituito.

\subsection{Che cos'è davvero una classe}

Non è una nuova entità del linguaggio: \code{typeof(Pet)} restituisce
\code{"function"}. In JavaScript le classi sono un tipo particolare di
\textbf{funzione}.

Che cosa fa allora il costrutto \code{class Pet \{...\}}?

\begin{enumerate}
  \item Crea una funzione costruttore chiamata \code{Pet}, il cui codice è preso
        dal metodo \code{constructor()} (o è vuoto se il metodo non c'è).
  \item Memorizza gli altri metodi della classe in \code{Pet.prototype}.
\end{enumerate}

È esattamente il secondo dei due modi di aggiungere metodi visti all'inizio del
capitolo: le classi sono \textbf{zucchero sintattico} sopra i prototipi.

\subsection{Non solo zucchero sintattico}

Il risultato è molto simile a quello di una funzione costruttore scritta a mano,
ma non identico. Le differenze principali:

\begin{itemize}
  \item una classe \textbf{non si può chiamare senza} \code{new}: si ottiene un
        \code{TypeError}, mentre una funzione costruttore chiamata per errore
        senza \code{new} viene eseguita come funzione normale (il problema visto
        con \code{chuck.constructor('Bob', 1)});
  \item i metodi della classe \textbf{non sono enumerabili}: non compaiono in
        un \code{for..in} sull'oggetto, a differenza di quelli aggiunti con
        \code{Pet.prototype.metodo = ...};
  \item il corpo della classe è sempre in modalità \textbf{strict}, anche se lo
        script non lo è;
  \item una classe non si può usare \textbf{prima} della sua dichiarazione,
        come le variabili \code{let}.
\end{itemize}

\begin{lstlisting}[style=js, name=classe-vs-funzione.js]
class Dog {
  constructor(name) { this.name = name; }
  bark() { return "Woof"; }
}

Dog("Rex");   // TypeError: Class constructor Dog cannot be
              //            invoked without 'new'

let rex = new Dog("Rex");
for (let key in rex) console.log(key);   // solo "name", non "bark"
\end{lstlisting}

\subsection{Proprietà, getter e setter}

Come gli oggetti, anche le classi supportano proprietà e accessori.

\begin{affianco}[0.5][c]
\begin{lstlisting}[style=js, name=classi-props.js]
class Duck {
  species = "Duck";                    // proprieta'

  constructor(name) { this.name = name }

  greet() {
    console.log("Quack!");
  }

  get fullDescription() {
    return `${this.name} the ${this.species}`;
  }

  [Symbol.iterator]() { /* ... */ }
}

let duck = new Duck("Chuck");
duck.greet();                        // Quack!
console.log(duck.fullDescription);   // Chuck the Duck
\end{lstlisting}
\risultato
\begin{immagine}[\code{species} e \code{name} sono proprietà di ogni\\
                 oggetto; i metodi e il getter stanno sul prototipo.]
\begin{tikzpicture}[
    pr/.style={-{Stealth[length=5pt,width=4pt]}, line width=0.85pt,
      draw=accent!70!black}]
  \oggetto[fill=accent!14!white]{pp}{(0,0)}{3.2cm}{Duck.prototype}
    {greet/function, fullDescription/get, [Symbol.iterator]/function}
  \oggetto{d}{(0,-2.6)}{3.2cm}{duck}{species/"Duck", name/"Chuck"}
  \draw[pr] (d.north) -- node[wtlab, right, xshift=2pt]{\code{[[Prototype]]}}
    (pp-box.south);
\end{tikzpicture}
\end{immagine}
\end{affianco}

\subsection{Ereditarietà fra classi}

Le classi JavaScript possono ereditare da altre classi con la parola chiave
\code{extends}.

Internamente, \code{extends} è realizzato proprio tramite la proprietà
\code{[[Prototype]]}: il prototipo di \code{Duck.prototype} è
\code{Pet.prototype}.

\begin{affianco}[0.5][c]
\begin{lstlisting}[style=js, name=ereditarieta.js]
class Pet {
  describe() { console.log("I'm a pet"); }
  greet()    { console.log("..."); }
}

class Duck extends Pet {
  greet() { console.log("Quack!"); }   // sovrascrive
}

let chuck = new Duck();
chuck.describe();   // I'm a pet   (ereditato)
chuck.greet();      // Quack!      (sovrascritto)
\end{lstlisting}
\risultato
\begin{immagine}[L'ereditarietà fra classi è ereditarietà fra prototipi.]
\begin{tikzpicture}[
    pr/.style={-{Stealth[length=5pt,width=4pt]}, line width=0.85pt,
      draw=accent!70!black}]
  \mchiave{1.4cm}
  \oggetto{c}{(0,0)}{2.6cm}{chuck}{}
  \oggetto[fill=accent!14!white]{dp}{(0,-1.55)}{2.6cm}{Duck.prototype}
    {greet/function}
  \oggetto[fill=accent!14!white]{pp}{(0,-3.35)}{2.6cm}{Pet.prototype}
    {describe/function, greet/function}
  \oggetto[fill=neutral!12!white]{op}{(0,-5.55)}{2.6cm}{Object.prototype}{}
  \draw[pr] (c-box.south) -- node[wtlab, right, xshift=2pt]{\code{[[Prototype]]}} (dp.north);
  \draw[pr] (dp-box.south) -- node[wtlab, right, xshift=2pt]{\code{[[Prototype]]}} (pp.north);
  \draw[pr] (pp-box.south) -- node[wtlab, right, xshift=2pt]{\code{[[Prototype]]}} (op.north);
  \draw[mcerca] ([xshift=-4pt]c-1.west) to[out=180, in=180]
    node[mnota, left]{\code{greet}?} ([xshift=-4pt]dp-1.west);
  \draw[mcerca] ([xshift=-4pt]c-1.west) to[out=190, in=180, looseness=1.4]
    node[mnota, left, pos=0.6]{\code{describe}?} ([xshift=-4pt]pp-1.west);
  \node[mnota, anchor=west, text=neutral!45!black] at ([xshift=4pt]pp-2.east)
    {(oscurato)};
\end{tikzpicture}
\end{immagine}
\end{affianco}

\section{Gestione degli errori}

Quando uno script viene eseguito, possono verificarsi errori: perché chi
programma sbaglia, perché l'utente fornisce input inattesi, e così via. Quando
accade, di norma lo script si ferma immediatamente e stampa un messaggio in
console.

\begin{affianco}[0.55][c]
\begin{lstlisting}[style=js, name=errore.js]
let myvar = 0;
for (let num of [1, 2, 3, 4, 5]) {
  mvvar += num;    // refuso: mvvar non esiste!
}
console.log(myvar);
\end{lstlisting}
\risultato
\begin{console}
Uncaught ReferenceError: mvvar is not defined
\end{console}
\wtnota{Lo script si ferma: \code{console.log(myvar)}\\ non viene mai eseguita.}
\end{affianco}

\subsection{try / catch / finally}

Il costrutto \code{try/catch/finally} permette di gestire gli errori.
Concettualmente è lo stesso costrutto di Java, con qualche differenza.

\begin{affianco}[0.55][c]
\begin{lstlisting}[style=js, name=trycatch.js]
"use strict";
try {
  console.log("starting");
  x = 1;                            // manca il let
  console.log("assignment done");   // non viene eseguita
} catch (error) {
  console.log(error.name);
  console.log(error.message);
} finally {
  console.log("done");
}
\end{lstlisting}
\risultato
\begin{immagine}
\begin{tikzpicture}[
    blk/.style={rounded corners=3pt, draw=primary!65!black, fill=primary!6!white,
      line width=0.7pt, font=\FiraCodeNL\scriptsize, minimum width=2.1cm,
      minimum height=0.6cm},
    a/.style={-{Stealth[length=5pt,width=4pt]}, line width=0.85pt,
      draw=merr},
    b/.style={-{Stealth[length=5pt,width=4pt]}, line width=0.8pt,
      draw=mok, dashed}]
  \node[blk] (t) at (0,0) {try};
  \node[blk, draw=merr, fill=red-gradient-1!35!white] (c) at (0,-1.3) {catch (error)};
  \node[blk] (f) at (0,-2.6) {finally};
  \draw[a] (t) -- node[mnota, left, text=merr]{errore} (c);
  \draw[a] (c) -- (f);
  \draw[b] (t.east) to[out=-30, in=30] node[mnota, right, text=mok, align=left]
    {nessun\\ errore} (f.east);
  \node[mnota, anchor=east] at ([xshift=-4pt]t.west) {\code{x = 1} \merr};
\end{tikzpicture}
\end{immagine}
\begin{console}
starting
ReferenceError
x is not defined
done
\end{console}
\end{affianco}



\warning{L'errore c'è solo in modalità strict}{
  Le slide omettono la riga \code{"use strict"}, che qui è stata aggiunta. Senza
  di essa, assegnare una variabile mai dichiarata \textbf{non} è un errore:
  JavaScript crea in silenzio una variabile globale \code{x}, stampa
  \textit{assignment done}, salta il \code{catch} ed esegue il \code{finally}.
  Lo stesso vale nella console dei DevTools, che non è in modalità strict. Il
  comportamento delle slide si ottiene invece senza la direttiva dentro un
  modulo o il corpo di una classe, che sono sempre strict.}

Le particolarità rispetto a Java sono due:

\begin{itemize}
  \item ci può essere \textbf{al più un} blocco \code{catch} (è comunque
        ammesso \code{try/finally} senza \code{catch});
  \item i diversi tipi di errore si distinguono \textbf{dentro} il blocco
        \code{catch}, non con più blocchi.
\end{itemize}

\begin{lstlisting}[style=js, name=tipi-errore.js]
"use strict";
try {
  x = 1;
} catch (error) {
  if (error instanceof ReferenceError) {
    console.log("Got a ReferenceError");
  } else {
    console.log(`Got something else: ${error.name}`);
  }
}
\end{lstlisting}

\subsection{L'operatore instanceof}

L'operatore \code{instanceof} verifica se la proprietà \code{prototype} di un
dato costruttore compare \textbf{da qualche parte} nella catena dei prototipi di
un oggetto.

\begin{affianco}[0.45][c]
\begin{lstlisting}[style=js, name=instanceof.js]
function Pet(name, species, age) {
  this.name = name;
  this.species = species;
  this.age = age;
}

const pet = new Pet('Chuck', 'Duck', 3);

console.log(pet instanceof Pet);      // true
console.log(pet instanceof Object);   // true
\end{lstlisting}
\risultato
\begin{immagine}[\code{pet instanceof Pet} e \code{pet instanceof Object}:\\
                 risalendo la catena (in viola) si incontrano\\
                 sia \code{Pet.prototype} sia \code{Object.prototype}.]
\begin{tikzpicture}[font=\Lato\footnotesize,
      o/.style={rounded corners=3pt, draw=primary!65!black, line width=0.8pt,
                fill=primary!8!white, minimum width=3.1cm, minimum height=1.0cm,
                align=center},
      pt/.style={o, draw=accent!60!black, fill=accent!9!white},
      g/.style={o, draw=neutral!55!white, fill=neutral!7!white},
      f/.style={o, minimum width=2.2cm},
      a/.style={-{Stealth[length=5pt,width=4pt]}, line width=0.85pt,
                draw=neutral!45!black},
      ch/.style={a, draw=accent!70!black}]

    \node[o]  (pe) at (0,0)    {\code{pet}\\\scriptsize name, species, age};
    \node[pt] (pp) at (0,1.9)  {\code{Pet.prototype}\\\scriptsize constructor: Pet};
    \node[g]  (op) at (0,3.8)  {\code{Object.prototype}\\\scriptsize toString(), valueOf()\ldots};

    \node[f] (fp) at (5.6,1.9) {\code{Pet}};
    \node[f] (fo) at (5.6,3.8) {\code{Object}};

    \draw[ch] (pe) -- node[wtlab, left, xshift=-2pt]{\code{[[Prototype]]}} (pp);
    \draw[ch] (pp) -- node[wtlab, left, xshift=-2pt]{\code{[[Prototype]]}} (op);

    \draw[a] ([yshift=6pt]fp.west) -- node[wtlab, above]{\code{prototype}}
             ([yshift=6pt]pp.east);
    \draw[a] ([yshift=-6pt]pp.east) -- node[wtlab, below]{\code{constructor}}
             ([yshift=-6pt]fp.west);
    \draw[a] ([yshift=6pt]fo.west) -- node[wtlab, above]{\code{prototype}}
             ([yshift=6pt]op.east);
    \draw[a] ([yshift=-6pt]op.east) -- node[wtlab, below]{\code{constructor}}
             ([yshift=-6pt]fo.west);
  \end{tikzpicture}
\end{immagine}
\end{affianco}

Perché anche \code{pet instanceof Object} è vero? Basta seguire la catena dei
prototipi di \code{pet}, a destra. Il primo anello è \code{Pet.prototype}, cioè il
prototipo predefinito di \code{Pet}; il secondo è \code{Object.prototype}.
Entrambi i costruttori hanno il proprio \code{prototype} nella catena.


In altre parole, \code{obj instanceof C} equivale a questa funzione, che risale
la catena un anello alla volta finché non trova \code{C.prototype} o arriva in
fondo (\code{null}):

\begin{lstlisting}[style=js, name=instanceof-a-mano.js]
function myInstanceof(obj, C) {
  let p = Object.getPrototypeOf(obj);
  while (p !== null) {
    if (p === C.prototype) return true;
    p = Object.getPrototypeOf(p);
  }
  return false;
}

myInstanceof(pet, Pet);      // true
myInstanceof(pet, Object);   // true
myInstanceof(pet, Array);    // false
\end{lstlisting}

È lo stesso meccanismo che fa funzionare il controllo nel \code{catch}: nella
catena di un \code{ReferenceError} ci sono sia il suo prototipo sia quello di
\code{Error}, quindi sono veri entrambi i controlli.

\begin{lstlisting}[style=js]
error instanceof ReferenceError;   // true
error instanceof Error;            // true: ogni errore "eredita" da Error
\end{lstlisting}

\subsection{Sollevare errori}

Gli errori in JavaScript si propagano come in Java: verso l'alto, finché non
vengono catturati o finché non si raggiunge la cima dell'albero delle chiamate
(nel qual caso l'esecuzione dello script viene interrotta bruscamente e l'errore
finisce in console). Si sollevano con la parola chiave \code{throw}.

\begin{affianco}[0.5][c]
\begin{lstlisting}[style=js, name=throw.js]
try {
  throw new Error("Oops!");
} catch (err) {
  console.log(err.name);      // Error
  console.log(err.message);   // Oops!
}
\end{lstlisting}
\risultato
\begin{immagine}[\code{throw} lancia un oggetto: \code{catch} lo riceve\\
                 nel suo parametro.]
\begin{tikzpicture}[
    a/.style={-{Stealth[length=5pt,width=4pt]}, line width=0.85pt, draw=merr}]
  \mchiave{1.2cm}
  \node[mcod] (t) at (0,0) {throw new Error("Oops!")};
  \oggetto[fill=red-gradient-1!45!white, draw=merr]{e}{(0,-1.25)}{2.6cm}{Error}
    {name/"Error", message/"Oops!"}
  \draw[a] (t) -- (e.north);
  \variabile[minimum width=0.55cm]{v}{(-2.5,-1.48)}{err}{}
  \riferimento{v}{-- (e-box.west |- v)}
\end{tikzpicture}
\end{immagine}
\end{affianco}

È possibile definire tipi di errore propri, estendendo \code{Error}:

\begin{affianco}[0.5][c]
\begin{lstlisting}[style=js, name=errore-custom.js]
class MyError extends Error {
  constructor(message = "Whoops") {
    super(message);
    this.name = "MyError";
  }
}

try {
  throw new MyError();
} catch (err) {
  console.log(err.name);      // MyError
  console.log(err.message);   // Whoops
}
\end{lstlisting}
\risultato
\begin{immagine}[\code{err} è sia un \code{MyError} sia un \code{Error}:\\
                 entrambi i prototipi sono nella sua catena.]
\begin{tikzpicture}[
    pr/.style={-{Stealth[length=5pt,width=4pt]}, line width=0.85pt,
      draw=accent!70!black}]
  \mchiave{1.2cm}
  \oggetto[fill=red-gradient-1!45!white]{e}{(0,0)}{2.6cm}{err}
    {name/"MyError", message/"Whoops"}
  \oggetto[fill=accent!14!white]{mp}{(0,-1.85)}{2.6cm}{MyError.prototype}{}
  \oggetto[fill=accent!14!white]{ep}{(0,-3.35)}{2.6cm}{Error.prototype}{}
  \oggetto[fill=neutral!12!white]{op}{(0,-4.85)}{2.6cm}{Object.prototype}{}
  \draw[pr] (e-box.south) -- (mp.north);
  \draw[pr] (mp-box.south) -- (ep.north);
  \draw[pr] (ep-box.south) -- (op.north);
  \node[mnota, anchor=west] at ([xshift=6pt]mp-box.east) {\code{instanceof MyError}};
  \node[mnota, anchor=west] at ([xshift=6pt]ep-box.east) {\code{instanceof Error}};
\end{tikzpicture}
\end{immagine}
\end{affianco}

\section{Moduli}

\subsection{Perché servono}

Il JavaScript moderno permette di definire \term{moduli} a livello di
linguaggio: sono un modo di suddividere programmi complessi in più file,
ciascuno contenente classi o funzioni per uno scopo specifico.

I vantaggi:

\begin{itemize}
  \item migliorano la \textbf{manutenibilità};
  \item favoriscono la \textbf{separazione delle responsabilità};
  \item aiutano a disambiguare e a prevenire i \textbf{conflitti di nome},
        migliorando la riusabilità.
\end{itemize}

Un modulo è semplicemente un file JavaScript. I moduli si caricano a vicenda e
usano le direttive \code{export} e \code{import} per scambiarsi funzionalità.

\begin{itemize}
  \item \code{export} marca variabili e funzioni che devono essere accessibili
        fuori dal modulo corrente;
  \item \code{import} permette di importare funzionalità specifiche da altri
        moduli, purché siano state esportate.
\end{itemize}

\subsection{Il problema, e la soluzione}

\begin{esempio}{Il conflitto di nomi}
Supponiamo di voler riutilizzare del codice scritto per altri progetti: un file
\code{pet.js} e un file \code{greet.js}.

\begin{affianco}[0.52]
\begin{lstlisting}[style=js, name=pet.js]
function Pet(name, age) {
  this.name = name; this.age = age;
}

let msg = "Hello";

function greet(pet) {
  console.log(`${msg}, I'm ${pet.name}`);
}
\end{lstlisting}
\risultato
\begin{lstlisting}[style=js, name=greet.js]
let msg = "Howdy";

function greet(name, message = msg) {
  console.log(`${message}, ${name}`);
}
\end{lstlisting}
\end{affianco}

\begin{lstlisting}[style=js, name=script.js]
let chuck = new Pet("Chuck", 7);   // Pet e' definita in pet.js
greet(chuck);                      // vorremmo: Hello, I'm Chuck
greet("Web Technologies");         // vorremmo: Howdy, Web Technologies
\end{lstlisting}

Si potrebbe pensare di includere semplicemente tutti gli script nella pagina:

\begin{affianco}[0.45][c]
\begin{lstlisting}[style=html]
<script src="./pet.js"></script>
<script src="./greet.js"></script>
<script src="./script.js"></script>
\end{lstlisting}
\risultato
\begin{immagine}[Un solo scope globale per tutti gli script:\\
                 \code{greet.js} sovrascrive \code{msg} e \code{greet}.]
\begin{tikzpicture}[
    r/.style={mcod, anchor=west, minimum height=0.48cm},
    da/.style={mnota, anchor=west, text=neutral!45!black}]
  \node[r] (r1) at (0,0)     {Pet};
  \node[da] at (2.6,0)       {da \code{pet.js}};
  \node[r] (r2) at (0,-0.55) {msg = "Howdy"};
  \node[da, text=merr] at (2.6,-0.55) {il \code{"Hello"} di \code{pet.js} è perso};
  \node[r] (r3) at (0,-1.1)  {greet(name, ...)};
  \node[da, text=merr] (n3) at (2.6,-1.1) {la \code{greet} di \code{pet.js} è persa};
  \scopo{1}{g}{scope globale}{(r1)(r2)(r3)(n3)}
\end{tikzpicture}
\end{immagine}
\end{affianco}

Ma non funziona: entrambi i file definiscono \code{msg} e \code{greet} nello
stesso scope globale, e il secondo sovrascrive il primo. Per farlo funzionare
dovremmo \textbf{modificare} il codice che volevamo riutilizzare, usando
identificatori diversi. Il che vanifica lo scopo stesso del riutilizzo.
\end{esempio}

\begin{esempio}{La soluzione con i moduli}
\begin{affianco}[0.52]
\begin{lstlisting}[style=js, name=pet-module.js]
export function Pet(name, age) {
  this.name = name; this.age = age;
}

let msg = "Hello";   // non esportata

export function greet(pet) {
  console.log(`${msg}, I'm ${pet.name}`);
}
\end{lstlisting}
\risultato
\begin{lstlisting}[style=js, name=greet-module.js]
let msg = "Howdy";   // non esportata

export function greet(name, message = msg) {
  console.log(`${message}, ${name}`);
}
\end{lstlisting}
\end{affianco}

Nel documento HTML basta includere lo script principale, dichiarandolo come
modulo:

\begin{lstlisting}[style=html]
<script type="module" src="./script-module.js"></script>
\end{lstlisting}

\begin{affianco}[0.52][c]
\begin{lstlisting}[style=js, name=script-module.js]
import { Pet, greet as greetPet } from './pet-module.js';
import { greet } from './greet-module.js';

let chuck = new Pet("Chuck", 7);
greetPet(chuck);              // Hello, I'm Chuck
greet("Web Technologies");    // Howdy, Web Technologies
\end{lstlisting}
\risultato
\begin{immagine}[Ogni modulo ha il suo scope: \code{msg} resta privato,\\
                 passa solo ciò che è esportato e importato.]
\begin{tikzpicture}[
    r/.style={mcod, anchor=west, minimum height=0.45cm},
    ex/.style={-{Stealth[length=5pt,width=4pt]}, line width=0.8pt,
      draw=accent!70!black}]
  % pet-module.js
  \node[r] (p1) at (0,0) {export Pet};
  \node[r] (p2) at (0,-0.5) {export greet};
  \node[r, text=neutral!50!black] (p3) at (0,-1.0) {msg = "Hello"};
  \scopo{2}{pm}{pet-module.js}{(p1)(p2)(p3)}
  % greet-module.js
  \node[r] (g1) at (0,-2.3) {export greet};
  \node[r, text=neutral!50!black] (g2) at (0,-2.8) {msg = "Howdy"};
  \scopo{2}{gm}{greet-module.js}{(g1)(g2)}
  % script-module.js
  \node[r] (s1) at (4.0,-0.6) {Pet};
  \node[r] (s2) at (4.0,-1.1) {greetPet};
  \node[r] (s3) at (4.0,-1.6) {greet};
  \scopo{3}{sm}{script-module.js}{(s1)(s2)(s3)}
  \draw[ex] (p1.east) to[out=0, in=180] (s1.west);
  \draw[ex] (p2.east) to[out=0, in=180] node[mnota, above, sloped, pos=0.55]{as} (s2.west);
  \draw[ex] (g1.east) to[out=0, in=180] (s3.west);
\end{tikzpicture}
\end{immagine}
\end{affianco}

Due cose da notare. Le variabili \code{msg} non danno più conflitto, perché
ciascuna vive nello \textbf{scope del proprio modulo}. E la clausola \code{as}
permette di rinominare un import, risolvendo il conflitto fra le due funzioni
\code{greet} senza toccare i file originali.
\end{esempio}

\info{I moduli hanno un proprio scope}{
  È questo il punto centrale: ogni modulo ha uno scope suo. Ciò che non viene
  esportato non è visibile dall'esterno, e ciò che viene importato va indicato
  esplicitamente. Il vecchio modello, in cui tutti gli script condividevano un
  unico scope globale, rendeva impossibile riusare codice scritto da altri senza
  leggerlo per intero a caccia di collisioni.}

\subsection{Altre forme di export e import}

Oltre a marcare con \code{export} ogni dichiarazione, si possono esportare più
nomi insieme in fondo al file, e ogni modulo può avere al più un
\term{export predefinito} (\code{export default}), che si importa senza
parentesi graffe e con il nome che si preferisce.

\begin{affianco}[0.45]
\begin{lstlisting}[style=js, name=geometria.js]
const PI = 3.14;
function area(r) { return PI * r * r; }

export { PI, area };       // piu' nomi insieme
export default class Circle { // predefinito
  /* ... */
}
\end{lstlisting}
\risultato
\begin{lstlisting}[style=js, name=main.js]
import Circle, { area } from './geometria.js';
                            // default + nominato
import * as geo from './geometria.js';
                            // tutto in un oggetto

area(1);    // 3.14
geo.PI;     // 3.14
\end{lstlisting}
\end{affianco}

\subsection{Cose da sapere sui moduli}

\begin{itemize}
  \item Un modulo è sempre in modalità \textbf{strict}, anche senza
        \code{"use strict"}; al livello più esterno \code{this} vale
        \code{undefined}.
  \item Uno script di tipo modulo è \textbf{differito}: viene eseguito solo
        dopo che il documento è stato analizzato, come con l'attributo
        \attr{defer}.
  \item Un modulo viene eseguito \textbf{una volta sola}, anche se è importato
        da più file: tutti ricevono le stesse funzioni e variabili.
  \item Gli import sono in sola lettura: assegnare un valore a \code{greetPet}
        nello script solleva un \code{TypeError}.
  \item Il percorso dopo \code{from} è un URL, relativo al file che importa:
        va scritto per intero, estensione \code{.js} compresa.
\end{itemize}

\warning{I moduli non funzionano aprendo il file con un doppio clic}{
  Se si apre la pagina direttamente dal disco (indirizzo \code{file://...}),
  il browser blocca il caricamento dei moduli per motivi di sicurezza e in
  console compare un errore CORS. La pagina va servita da un server web,
  anche solo locale: per esempio con l'estensione \textit{Live Server} di
  VS Code, oppure con \code{npx serve} o \code{python3 -m http.server} lanciati
  nella cartella del progetto.}

\section{Riferimenti}

\begin{itemize}
  \item \textbf{The Modern JavaScript Tutorial} \\
        \urlx{https://javascript.info/} \\
        Parte 1: \textit{Prototypes and inheritance}, \textit{Data types}
        (5.4--5.7), \textit{Classes} (9.1--9.4, 9.6), \textit{Generators and
        advanced iteration} (12.1), \textit{Modules} (13.1, 13.2).
  \item \textbf{Eloquent JavaScript} (3\textsuperscript{a} edizione), Marijn
        Haverbeke. \\
        \urlx{https://eloquentjavascript.net/} \\
        Capitoli 6 e 10.
  \item \textbf{JavaScript Reference} --- MDN web docs. \\
        \urlx{https://developer.mozilla.org/en-US/docs/Web/JavaScript/Reference}
\end{itemize}
